%============================================================
% CH1 · P2 · PROBLEM 1 — FIND IF PATH EXISTS IN GRAPH (LeetCode 1971)
%============================================================
\subsection{Problem 1.2.1: Find if Path Exists in Graph}
\problemlabel{1.2.1}{ssec:ch1-path-exists}

\problemcard{1.2.1}{Find if Path Exists in Graph}{LeetCode 1971}{https://leetcode.com/problems/find-if-path-exists-in-graph/}{\exeasy\ \lc{Easy}\ (the canonical reachability question)}{One BFS/DFS from the source with early exit; DSU as an alternative}

\subsubsection*{1. Problem statement}

An undirected graph has $n$ vertices labelled $0..n-1$ and an edge list \texttt{edges}, where $\texttt{edges}[i] = [u_i, v_i]$ is a bi-directional edge. Return \texttt{true} if there is a valid path from vertex \texttt{source} to vertex \texttt{destination}, and \texttt{false} otherwise.

\textbf{Constraints:} $1 \le n \le 2\cdot 10^5$, \; $0 \le |\texttt{edges}| \le 2\cdot 10^5$, \; $0 \le u_i, v_i \le n-1$, \; $u_i \ne v_i$ (no self-loops), no duplicate edges, $0 \le \texttt{source}, \texttt{destination} \le n-1$.

\begin{table}[H]
\centering\small
\begin{tabular}{@{}clccl@{}}
\toprule
$n$ & \textbf{edges} & \textbf{source} & \textbf{destination} & \textbf{Output} \\
\midrule
3 & $[[0,1],[1,2],[2,0]]$ & 0 & 2 & \texttt{true} \quad ($0 \to 2$ directly, or $0 \to 1 \to 2$) \\
6 & $[[0,1],[0,2],[3,5],[5,4],[4,3]]$ & 0 & 5 & \texttt{false} \quad (5 is in another piece) \\
\bottomrule
\end{tabular}
\end{table}

\begin{diagrambox}
\begin{tikzpicture}
    \begin{scope}
        \node[gvsrc] (a0) at (0,0) {0};
        \node[gv] (a1) at (1.4,0.8) {1};
        \node[gvdst] (a2) at (1.4,-0.8) {2};
        \draw[ge] (a0) -- (a1); \draw[ge] (a1) -- (a2); \draw[ge] (a2) -- (a0);
        \node[note] at (0.7,-1.6) {Sample 1: source 0, destination 2\\answer \texttt{true}};
    \end{scope}
    \begin{scope}[xshift=5cm]
        \node[gvsrc] (b0) at (0,0) {0};
        \node[gv] (b1) at (1.4,0.8) {1};
        \node[gv] (b2) at (1.4,-0.8) {2};
        \node[gv] (b3) at (3.6,0.8) {3};
        \node[gvdst] (b5) at (5,0) {5};
        \node[gv] (b4) at (3.6,-0.8) {4};
        \draw[ge] (b0) -- (b1); \draw[ge] (b0) -- (b2);
        \draw[ge] (b3) -- (b5); \draw[ge] (b5) -- (b4); \draw[ge] (b4) -- (b3);
        \draw[gcut] (2.5,1.3) -- (2.5,-1.3);
        \node[note] at (2.5,-1.6) {Sample 2: source 0, destination 5\\no edge crosses the red line $\Rightarrow$ \texttt{false}};
    \end{scope}
\end{tikzpicture}
\figcap{The two official samples. In Sample~2 the set $\{0,1,2\}$ has no edge leaving it, which is a complete proof that 5 is unreachable.}
\end{diagrambox}

\subsubsection*{2. The graph model}

\begin{modelbox}[Graph Model --- Find if Path Exists in Graph]
\renewcommand{\arraystretch}{1.35}
\begin{tabularx}{\linewidth}{@{}>{\bfseries}l X@{}}
Vertices & the $n$ labelled vertices $0..n-1$ --- the graph is given explicitly \\
Edges & every $[u, v]$ is undirected: store it as $u \to v$ \emph{and} $v \to u$ (adjacency lists) \\
Weights & none \\
Question & is \texttt{destination} $\in R(\texttt{source})$? Equivalently: are the two vertices in the same connected component? \\
Size & $|V| \le 2\cdot10^5$, $|E| \le 2\cdot 10^5$ ($4\cdot10^5$ adjacency entries): any $\VE$ algorithm runs in a few milliseconds; $\bigO{|V|\cdot|E|} = 4\cdot10^{10}$ is far too slow \\
\end{tabularx}
\end{modelbox}

The input is an \emph{edge list}; a search needs \emph{adjacency lists} (``who are the neighbours of $u$?''). Converting one into the other is the first five lines of every solution:

\begin{diagrambox}
\begin{tikzpicture}[lst/.style={font=\small\ttfamily, anchor=west}]
    \node[font=\small\bfseries, anchor=west] at (0,1.9) {edge list (input)};
    \foreach \e [count=\i] in {{[0,1]},{[0,2]},{[3,5]},{[5,4]},{[4,3]}}
        \node[lst] at (0,1.9-0.45*\i) {\e};
    \draw[flow] (1.6,0.6) -- node[lbl, above] {each edge twice} (4.0,0.6);
    \node[font=\small\bfseries, anchor=west] at (4.3,1.9) {adjacency lists (built)};
    \foreach \v/\l [count=\i] in {0/{1, 2}, 1/{0}, 2/{0}, 3/{5, 4}, 4/{5, 3}, 5/{3, 4}}
        \node[lst] at (4.3,1.9-0.4*\i) {adj[\v] = \{\l\}};
\end{tikzpicture}
\figcap{Sample 2. Edge $[5,4]$ produces two entries: \texttt{adj[5]} gets 4 and \texttt{adj[4]} gets 5. The order inside each list follows the input order.}
\end{diagrambox}

\subsubsection*{3. Key observations}

\begin{itemize}[leftmargin=*]
    \item \textbf{Only yes/no is asked.} We do not need the path, its length or the number of paths --- exactly the reach-set question of the pattern.
    \item \textbf{Local information is not enough.} ``Both endpoints have at least one edge'' is true in Sample~2, yet the answer is \texttt{false}. Reachability depends on whole chains of edges.
    \item \textbf{$\texttt{source} = \texttt{destination}$ is \texttt{true}}, even with no edges at all (a walk of zero edges). With $n = 1$ this is the only possible query.
    \item \textbf{Early exit.} The search may stop the moment \texttt{destination} is reached; on a \texttt{false} instance it must explore the whole component of \texttt{source}.
    \item \textbf{Depth.} A path graph $0 - 1 - \dots - (n{-}1)$ with $n = 2\cdot10^5$ is allowed, so a recursive DFS may go $2\cdot10^5$ calls deep. Prefer BFS or the iterative stack.
\end{itemize}

\subsubsection*{4. Approaches}

\begin{bruteforcebox}
\textbf{Idea: grow the reached set by sweeping the edge list until nothing changes.} Start with $\{\texttt{source}\}$. In one \emph{pass}, look at every edge $[u,v]$: if exactly one endpoint is reached, mark the other one. Repeat passes until a whole pass marks nothing.

\begin{lstlisting}[style=cpp]
#include <bits/stdc++.h>
using namespace std;

class Solution {
public:
    bool validPath(int n, vector<vector<int>>& edges, int source, int destination) {
        vector<char> reached(n, 0);
        reached[source] = 1;
        bool changed = true;
        while (changed) {                          // one pass over all edges
            changed = false;
            for (auto& e : edges) {
                int u = e[0], v = e[1];
                if (reached[u] && !reached[v]) { reached[v] = 1; changed = true; }
                else if (reached[v] && !reached[u]) { reached[u] = 1; changed = true; }
            }
        }
        return reached[destination];
    }
};
\end{lstlisting}

It is correct (when a pass changes nothing, no edge connects a reached and an unreached vertex --- the closure property of the pattern), but a pass may add only \emph{one} new vertex. On the path $0-1-2-3-4$ with the edges listed backwards, $[3,4], [2,3], [1,2], [0,1]$:

\begin{center}
\begin{tikzpicture}[x=1.05cm]
    \foreach \p/\k in {1/1, 2/2, 3/3, 4/4, 5/4}{
        \begin{scope}[yshift=-\p*0.62cm]
            \node[font=\scriptsize\bfseries, anchor=east] at (-0.5,0) {pass \p};
            \foreach \i in {0,...,4}{
                \ifnum\i>\k \node[gv, minimum size=5.5mm, font=\scriptsize\bfseries] (p\p\i) at (\i,0) {\i};
                \else \ifnum\i=0 \node[gvsrc, minimum size=5.5mm, font=\scriptsize\bfseries] (p\p\i) at (\i,0) {\i};
                \else \node[gvdone, minimum size=5.5mm, font=\scriptsize\bfseries] (p\p\i) at (\i,0) {\i}; \fi\fi
            }
            \foreach \i [evaluate=\i as \j using int(\i+1)] in {0,...,3} \draw[ge] (p\p\i) -- (p\p\j);
        \end{scope}
    }
    \node[note, anchor=west] at (4.6,-0.62) {edge $[0,1]$ is last in the list:\\only vertex 1 is added};
    \node[note, anchor=west] at (4.6,-3.1) {pass 5 changes nothing: stop};
\end{tikzpicture}
\end{center}

With $n$ vertices on such a path we need $n$ passes of $|E|$ edges each (the program counted exactly $1000$ passes for $n = 1000$).

\begin{complexitybox}
\textbf{Time:} $\bigO{|V| \cdot |E|}$ --- up to $2\cdot10^5 \times 2\cdot10^5 = 4\cdot10^{10}$ edge checks. \textbf{Space:} $\bigO{|V|}$. \textbf{Verdict:} \tle{}. (Enumerating all simple paths would be even worse: exponential.)
\end{complexitybox}
\end{bruteforcebox}

\begin{algorithmbox}[Approach 2: BFS from the source (early exit)]
\textbf{Idea.} Build adjacency lists once; then each vertex is processed once and each edge is looked at twice (once from each end). The brute force wasted time re-reading edges whose endpoints had not changed; a search only looks at the edges of a vertex at the moment that vertex becomes reached.

\begin{lstlisting}[style=cpp]
#include <bits/stdc++.h>
using namespace std;

class Solution {
public:
    bool validPath(int n, vector<vector<int>>& edges, int source, int destination) {
        vector<vector<int>> adj(n);
        for (auto& e : edges) {                    // undirected: store both directions
            adj[e[0]].push_back(e[1]);
            adj[e[1]].push_back(e[0]);
        }
        vector<char> seen(n, 0);
        queue<int> q;
        q.push(source);
        seen[source] = 1;
        while (!q.empty()) {
            int u = q.front(); q.pop();
            if (u == destination) return true;     // early exit (covers source == dest)
            for (int v : adj[u])
                if (!seen[v]) { seen[v] = 1; q.push(v); }
        }
        return false;                              // component of source fully explored
    }
};
\end{lstlisting}

\begin{complexitybox}
\textbf{Time:} $\bigO{|V| + |E|}$ --- building the lists is $\bigO{|V|+|E|}$, the search touches every vertex once and every adjacency entry at most once. \textbf{Space:} $\bigO{|V|+|E|}$ for the lists, $\bigO{|V|}$ for \texttt{seen} and the queue. On a path of $2\cdot 10^5$ vertices the program answered \texttt{true} in about $0.06$\,s, including building the graph. \textbf{Verdict:} \ac{}.
\end{complexitybox}
\end{algorithmbox}

\begin{invariantbox}[Why it is correct]
This is the pattern's search with one extra line, so its invariant applies: every vertex in the queue or already popped is reachable from \texttt{source} (soundness), and when the queue runs empty no edge leaves the marked set (closure), so the marked set is exactly the component of \texttt{source}.
\begin{itemize}[leftmargin=*]
    \item If the function returns \texttt{true}, \texttt{destination} was popped, hence marked, hence reachable.
    \item If it returns \texttt{false}, the queue emptied without popping \texttt{destination}. Every marked vertex is eventually popped, so \texttt{destination} was never marked, and by closure it is not reachable.
\end{itemize}
The early exit only stops the loop sooner; it never changes which answer is given.
\end{invariantbox}

\begin{alternativebox}[Approach 3: Disjoint Set Union (preview of Chapter 6)]
\textbf{Idea.} Forget searching: process the edges one by one and maintain the \emph{groups of vertices known to be connected}. Every vertex points to a \emph{parent}; following parents leads to the group's \emph{root} (its representative). An edge $[u,v]$ merges two groups by hanging one root under the other. At the end, \texttt{source} and \texttt{destination} are connected iff they have the same root.

\begin{lstlisting}[style=cpp]
#include <bits/stdc++.h>
using namespace std;

class Solution {
    vector<int> par;
    int findRoot(int x) {
        while (par[x] != x) {
            par[x] = par[par[x]];                  // path halving: shortcut every other step
            x = par[x];
        }
        return x;
    }
public:
    bool validPath(int n, vector<vector<int>>& edges, int source, int destination) {
        par.resize(n);
        for (int i = 0; i < n; i++) par[i] = i;   // every vertex is its own group
        for (auto& e : edges) {
            int a = findRoot(e[0]), b = findRoot(e[1]);
            if (a != b) par[a] = b;                 // merge the two groups
        }
        return findRoot(source) == findRoot(destination);
    }
};
\end{lstlisting}

On Sample~2 the parent array evolves as follows (output of the program):
\begin{center}\small
\begin{tabular}{@{}lcccccc@{}}
\toprule
\textbf{after edge} & \texttt{par[0]} & \texttt{par[1]} & \texttt{par[2]} & \texttt{par[3]} & \texttt{par[4]} & \texttt{par[5]} \\
\midrule
start     & 0 & 1 & 2 & 3 & 4 & 5 \\
$[0,1]$   & 1 & 1 & 2 & 3 & 4 & 5 \\
$[0,2]$   & 1 & 2 & 2 & 3 & 4 & 5 \\
$[3,5]$   & 1 & 2 & 2 & 5 & 4 & 5 \\
$[5,4]$   & 1 & 2 & 2 & 5 & 4 & 4 \\
$[4,3]$   & 1 & 2 & 2 & 4 & 4 & 4 \\
\bottomrule
\end{tabular}
\end{center}
At the last edge, \texttt{findRoot(3)} walks $3 \to 5 \to 4$ and shortcuts \texttt{par[3]} to $4$; both ends already share root $4$, so nothing is merged.

\begin{center}
\begin{tikzpicture}[x=1.2cm, y=0.95cm]
    \node[gvA] (r2) at (0,0) {2}; \node[gvA] (n1) at (0,-1) {1}; \node[gvsrc] (n0) at (0,-2) {0};
    \draw[gde] (n1) -- (r2); \draw[gde] (n0) -- (n1);
    \node[gvB] (r4) at (3,0) {4}; \node[gvB] (n3) at (2.4,-1) {3}; \node[gvdst] (n5) at (3.6,-1) {5};
    \draw[gde] (n3) -- (r4); \draw[gde] (n5) -- (r4);
    \node[note, anchor=west] at (4.3,-0.6) {root$(0) = 2 \ne 4 = $ root$(5)$\\$\Rightarrow$ \texttt{false}};
\end{tikzpicture}
\end{center}

\begin{complexitybox}
\textbf{Time:} $\bigO{(|V|+|E|)\,\log|V|}$ with path halving alone, practically linear; with union by size as well it is $\bigO{(|V|+|E|)\,\alpha(|V|)}$ (Chapter~6). \textbf{Space:} $\bigO{|V|}$ --- no adjacency lists at all. \textbf{Verdict:} \ac{}.
\end{complexitybox}

\textbf{When is DSU the better tool?} When edges arrive \emph{over time} mixed with ``are $a$ and $b$ connected now?'' queries, a search would have to be repeated after every new edge, while DSU just keeps merging. For a single query on a fixed graph, BFS and DSU are equally good.
\end{alternativebox}

\begin{tricksbox}[Approach 4: Implementation Refinements]
\textbf{Iterative DFS.} Replacing the queue by a \texttt{vector} used as a stack gives the same answer (the pattern's proof did not depend on the order):
\begin{lstlisting}[style=snippet]
vector<int> st(1, source);  seen[source] = 1;
while (!st.empty()) {
    int u = st.back(); st.pop_back();
    if (u == destination) return true;
    for (int v : adj[u]) if (!seen[v]) { seen[v] = 1; st.push_back(v); }
}
\end{lstlisting}
\textbf{Why not recursion here?} \texttt{dfs(u)} calling \texttt{dfs(v)} along the path $0 - 1 - \dots - 199999$ is $2 \cdot 10^5$ frames deep. Some judges survive it, others crash; the stack version costs nothing extra.

\textbf{Early exit at discovery time.} Testing \texttt{v == destination} when $v$ is \emph{marked} (instead of when popped) finishes one layer earlier. It then needs a separate check for $\texttt{source} = \texttt{destination}$; checking on pop, as in Approach~2, needs none.
\end{tricksbox}

\subsubsection*{5. Trace}

\begin{tracebox}[{Sample 2: $n = 6$, source $0$, destination $5$}]
Adjacency lists (input order): \texttt{adj[0]=\{1,2\}}, \texttt{adj[1]=\{0\}}, \texttt{adj[2]=\{0\}}, \texttt{adj[3]=\{5,4\}}, \texttt{adj[4]=\{5,3\}}, \texttt{adj[5]=\{3,4\}}.

\begin{center}
\begin{tikzpicture}[scale=0.62, transform shape]
    \foreach \k/\t/\a/\b/\c/\d/\e/\f in {
        0/{start: queue $[0]$}/gvsrc/gv/gv/gv/gv/gvdst,
        1/{pop 0, push 1, 2: queue $[1,2]$}/gvdone/gvseen/gvseen/gv/gv/gvdst,
        2/{pop 1, pop 2: queue empty}/gvdone/gvdone/gvdone/gvdead/gvdead/gvdead}{
        \begin{scope}[xshift=\k*8.3cm]
            \node[\a] (x0) at (0,0) {0};
            \node[\b] (x1) at (1.4,0.8) {1};
            \node[\c] (x2) at (1.4,-0.8) {2};
            \node[\d] (x3) at (3.6,0.8) {3};
            \node[\f] (x5) at (5,0) {5};
            \node[\e] (x4) at (3.6,-0.8) {4};
            \draw[ge] (x3) -- (x5); \draw[ge] (x5) -- (x4); \draw[ge] (x4) -- (x3);
            \ifnum\k>0 \draw[gtree] (x0) -- (x1); \draw[gtree] (x0) -- (x2);
            \else \draw[ge] (x0) -- (x1); \draw[ge] (x0) -- (x2); \fi
            \node[font=\large\bfseries] at (2.5,1.7) {\t};
        \end{scope}
    }
\end{tikzpicture}
\end{center}

\begin{center}\small
\begin{tabular}{@{}clll@{}}
\toprule
\textbf{step} & \textbf{pop} & \textbf{push (newly marked)} & \textbf{queue after (front on the left)} \\
\midrule
1 & 0 & 1, 2 & 1 2 \\
2 & 1 & --- \ (0 is marked) & 2 \\
3 & 2 & --- \ (0 is marked) & (empty) $\Rightarrow$ return \texttt{false} \\
\bottomrule
\end{tabular}
\end{center}
The search never looks at the edges $[3,5], [5,4], [4,3]$: it only pays for the component it explores. The finished set $\{0,1,2\}$ is the ``no edge leaves'' certificate drawn in the statement.
\end{tracebox}

\begin{tracebox}[{Sample 1: $n = 3$, source $0$, destination $2$}]
Adjacency lists: \texttt{adj[0]=\{1,2\}}, \texttt{adj[1]=\{0,2\}}, \texttt{adj[2]=\{1,0\}}.
\begin{center}
\begin{tikzpicture}[scale=0.8, transform shape]
    \foreach \k/\t/\a/\b/\c in {
        0/{start: queue $[0]$}/gvsrc/gv/gvdst,
        1/{pop 0: queue $[1,2]$}/gvdone/gvseen/gvseen,
        2/{pop 1: queue $[2]$}/gvdone/gvdone/gvseen,
        3/{pop 2 $=$ destination}/gvdone/gvdone/gvcur}{
        \begin{scope}[xshift=\k*4.1cm]
            \node[\a] (y0) at (0,0) {0};
            \node[\b] (y1) at (1.4,0.8) {1};
            \node[\c] (y2) at (1.4,-0.8) {2};
            \draw[ge] (y1) -- (y2);
            \ifnum\k>0 \draw[gtree] (y0) -- (y1); \draw[gtree] (y0) -- (y2);
            \else \draw[ge] (y0) -- (y1); \draw[ge] (y0) -- (y2); \fi
            \node[font=\small\bfseries] at (0.7,1.55) {\t};
        \end{scope}
    }
\end{tikzpicture}
\end{center}
Popping 1 finds both neighbours already marked. Popping 2 hits the early exit: \texttt{true}. The search tree ($0{-}1$, $0{-}2$) shows the path BFS found: $0 \to 2$, one edge.
\end{tracebox}

\subsubsection*{6. Test suite}

\begin{testbox}
\centering\small
\begin{tabular}{@{}p{6.6cm}cp{5.0cm}@{}}
\toprule
\textbf{Input} & \textbf{Expected} & \textbf{Why} \\
\midrule
$n=1$, no edges, $0 \to 0$ & \texttt{true} & zero-edge walk; smallest input \\
$n=2$, no edges, $0 \to 1$ & \texttt{false} & no edges at all \\
Sample 2 with source $5$, destination $3$ & \texttt{true} & the other component works \\
path $0-1-\dots-199999$, $0 \to 199999$ & \texttt{true} & depth $2\cdot10^5$: kills naive recursion \\
star with centre 0, $n = 2\cdot10^5$, $1 \to 2$ & \texttt{true} & one vertex with $2\cdot10^5-1$ neighbours \\
two paths of $10^5$ vertices each, $0 \to 199999$ & \texttt{false} & full exploration of a large component \\
\bottomrule
\end{tabular}
\end{testbox}
All six expected values were produced by both the BFS and the DSU program; in addition, BFS, iterative DFS, DSU and the brute force agreed with a Floyd--Warshall-style closure on 20\,000 random graphs with up to 8 vertices.

\subsubsection*{7. Common mistakes}

\begin{warningbox}[Mistake 1 --- one direction only]
\begin{lstlisting}[style=snippet]
for (auto& e : edges) adj[e[0]].push_back(e[1]);                          // WRONG
for (auto& e : edges) { adj[e[0]].push_back(e[1]); adj[e[1]].push_back(e[0]); }  // RIGHT
\end{lstlisting}
Both official samples still pass with the wrong version: Sample~1 keeps the cycle $0\to1\to2\to0$, and Sample~2 is \texttt{false} anyway. Sample~2's graph with source $1$ and destination $2$ catches it: the right answer is \texttt{true} ($1 - 0 - 2$), but the one-way lists give vertex 1 no neighbours at all, so the wrong version says \texttt{false} (checked by running both).
\end{warningbox}

\begin{warningbox}[Mistake 2 --- forgetting the trivial case]
Code that only checks \texttt{v == destination} for \emph{newly discovered} neighbours returns \texttt{false} for $\texttt{source} = \texttt{destination}$ (no neighbour of $0$ is $0$ itself; with $n=1$ there are no neighbours at all). Check when popping, or test \texttt{source == destination} first.
\end{warningbox}

\begin{warningbox}[Mistake 3 --- an adjacency matrix]
\texttt{vector<vector<bool>> g(n, vector<bool>(n))} with $n = 2\cdot 10^5$ needs $4\cdot 10^{10}$ bits $= 5$\,GB, and scanning a row costs $\bigO{n}$ per vertex, $\bigO{n^2}$ in total. Adjacency lists are $\bigO{|V|+|E|}$ in memory and time.
\end{warningbox}

\begin{warningbox}[Mistake 4 --- recursion on a long path]
A recursive DFS passes the samples and small random tests, then crashes on the $2\cdot10^5$-vertex path. Use BFS or an explicit stack (Approach~4).
\end{warningbox}

\subsubsection*{8. Takeaways}

\begin{itemize}[leftmargin=*]
    \item ``Is there a path from $a$ to $b$?'' $=$ ``is $b$ in $R(a)$?'' $=$ one search, $\VE$.
    \item Build adjacency lists from the edge list first; undirected $=$ both directions.
    \item Check the target when popping: it handles $\texttt{source} = \texttt{destination}$ for free.
    \item A failed search \emph{proves} impossibility: the explored set has no edge leaving it.
    \item DSU answers the same question without adjacency lists and shines when edges arrive over time (Chapter~6).
\end{itemize}

\begin{variantsbox}
\variant{One-way edges}{\varin{https://codeforces.com/problemset/problem/500/A}}
  {each edge can be used only from $u$ to $v$ (portals, one-way streets).}
  {same vertices; store only $u \to v$. Reachability is no longer symmetric.}
  {the same search; E2.1 is an instance where every vertex has a single out-edge, so the search becomes a loop.}
\variant{Many queries, fixed graph}{\varnone}
  {$q \le 2\cdot10^5$ questions ``is there a path from $a_i$ to $b_i$?'' on one undirected graph.}
  {label every vertex with the id of its connected component.}
  {one search per \emph{unlabelled} vertex, $\VE$ in total (Chapter~4); then each query is \texttt{comp[a] == comp[b]}, $\bigO{1}$.}
\variant{Edges arrive between queries}{\varnone}
  {the edge list is revealed one edge at a time, interleaved with connectivity questions.}
  {the same vertices; the graph only grows.}
  {DSU: \texttt{unite} for each edge, compare roots for each query, near-$\bigO{1}$ per operation (Chapter~6).}
\variant{Forbidden vertices}{\varnone}
  {the path may not pass through a given set $B$ of blocked vertices.}
  {delete $B$ from the graph --- or simply mark every vertex of $B$ as \texttt{seen} before the search.}
  {unchanged search; if \texttt{source} or \texttt{destination} is blocked, answer \texttt{false} (unless the statement says otherwise).}
\variant{Print a path}{\varnone}
  {output one path from \texttt{source} to \texttt{destination}, not just yes/no.}
  {store \texttt{parent[v] = u} at the moment $v$ is marked from $u$.}
  {walk \texttt{parent} back from \texttt{destination} and reverse; with BFS this path even has the fewest edges (Chapter~3).}
\end{variantsbox}
